\documentclass[11pt]{article}
\usepackage[T1]{fontenc}
\usepackage{lmodern}
\usepackage[margin=1in]{geometry}
\usepackage{amsmath,amssymb,amsthm,mathtools}
\usepackage[hidelinks]{hyperref}
\usepackage{microtype}
\numberwithin{equation}{section}
\newtheorem{theorem}{Theorem}[section]
\newtheorem{proposition}[theorem]{Proposition}
\newtheorem{lemma}[theorem]{Lemma}
\theoremstyle{remark}
\newtheorem{remark}[theorem]{Remark}
\newcommand{\R}{\mathbb R}
\newcommand{\C}{\mathbb C}
\newcommand{\PP}{\mathbb P}
\newcommand{\EE}{\mathbb E}
\newcommand{\Tr}{\operatorname{Tr}}
\newcommand{\divg}{\operatorname{div}}
\newcommand{\diag}{\operatorname{diag}}
\newcommand{\supp}{\operatorname{supp}}
\newcommand{\calE}{\mathcal E}
\newcommand{\calX}{\mathcal X}
\newcommand{\calM}{\mathcal M}
\newcommand{\norm}[1]{\left\lVert#1\right\rVert}
\newcommand{\ip}[2]{\left\langle#1,#2\right\rangle}
\DeclareMathOperator{\Rea}{Re}
\DeclareMathOperator{\argu}{arg}
\DeclareMathOperator{\Dom}{D}
\title{Local no-slip Navier--Stokes equations under stochastic parabolicity}
\author{Kaiwen Jin\\\small\href{mailto:yui.ui.siki@gmail.com}{yui.ui.siki@gmail.com}}
\date{19 October 2026}
\begin{document}
\maketitle

\begin{quote}\small
\textbf{Revision, 19 October 2026.}
The earlier version was recorded as partial progress after the
\href{https://github.com/MColbrook/AIM/blob/24edaadc6bb272737d8da0745c662c992db07b22/research/solution_reviews/2026-10-15-active-prs/458-review.md}{repository's independent AI review}
identified a missing hypothesis in Section~\ref{sec:channel-consumer}.
Lemmas~\ref{lem:channel-extra-derivative} and
\ref{lem:channel-compatible-right-inverses} now supply the compatible
fractional right-inverse estimates of the corrected comparison theorem
\cite[Theorem 1.3]{KWErratum}. This revised complete claim is submitted
for renewed independent review; the historical review is retained.
\end{quote}

\begin{abstract}
We consider the three-dimensional incompressible Navier--Stokes equation
on a bounded smooth domain with no-slip boundary condition and finitely
many smooth, divergence-free, boundary-tangent transport noises. We prove
local pathwise well-posedness under the stochastic parabolicity condition
$\nu I-\frac12\sum_k b_k\otimes b_k\geq\delta I>0$, without an
additional smallness assumption on the noise. For smooth compatible
initial data the solution is continuous in $B^{1/2}_{4,4}$ and belongs
locally to $L^4$ in time with values in $W^{1,4}_0$. Uniqueness holds
in this entire path class, and the solutions depend continuously in
probability on the initial data after common stopping. The proof first
controls constant transport on periodic channels and full interior
tori. Localization through the Stokes energy form preserves the weak
Dirichlet scale; a quantitative square-root comparison makes its errors
small enough to recover the physical-domain estimate.
\end{abstract}

\section{Introduction}\label{sec:intro}

Consider the It\^o equation
\begin{equation}\label{eq:NS}
 du=\bigl[\nu\PP_D\Delta u-\PP_D\divg(u\otimes u)\bigr]dt
 +\sum_{k=1}^K\PP_D[(b_k\cdot\nabla)u]\,dW^k_t,
 \qquad \divg u=0,\quad u|_{\partial D}=0.
\end{equation}
Here $D\subset\R^3$ is bounded and smooth, $\nu>0$ is the viscosity,
$\PP_D$ is the Helmholtz projection, and the finite family $b_k$ is
smooth, solenoidal, and tangent to the wall. Transport noise acts at
the same differential order as the square root of viscosity. Its
quadratic variation therefore changes the coercive second-order part
of an energy estimate. The natural positivity condition is
\[
 \nu I-\tfrac12\sum_kb_k(x)\otimes b_k(x)\geq\delta I>0
 \qquad(x\in\overline D).
\]
The question is whether this positivity suffices for a strong local
theory with no-slip walls, even when the transport cannot be treated
as a small perturbation.

The deterministic weak Stokes scale places convection in a negative
Sobolev space and reduces local existence to an estimate for the
linear equation. Pr\"uss and Wilke \cite[Section 5]{PW} develop this
formulation for deterministic Navier--Stokes equations. Stochastic
maximal regularity gives an analogous route once the linear noise
has been controlled; see Agresti and Veraar \cite{AV}. Their survey
distinguishes the whole-space non-small transport theory from the
no-slip boundary question, which appears as Open Problem 9. In the
presence of walls, pressure and the zero vector trace interact with
localization. In particular, Helmholtz projection enforces the normal
constraint but does not in general preserve the tangential zero trace.

Boundary results also depend on the precise equation and solution
class. Goodair's strong two-dimensional theory uses Navier boundary
conditions \cite{GoodairNavier}; its discussion explains the difficulty
of the corresponding no-slip estimates. His no-slip channel
inviscid-limit result \cite{GoodairPlates} uses weak martingale
solutions and a transport-stretching model. These are different
questions from pathwise local well-posedness for \eqref{eq:NS} on
arbitrary smooth three-dimensional domains. We address the finite
smooth-coefficient formulation in AIM Problem 458 \cite{AIM458}.
This is a specified part of the broader no-slip problem in \cite{AV}.

Let
\[
 \calE=\{h\in B^{1/2}_{4,4}(D;\R^3):\divg h=0,\ \Tr h=0\},
 \qquad V_4=W^{1,4}_{0,\sigma}(D),\qquad
 Y_4=(W^{1,4/3}_{0,\sigma}(D))^*.
\]
The trace in $\calE$ is well defined since $1/2>1/4$.
Equation \eqref{eq:NS} is understood as a weak solenoidal integral
identity in $Y_4$. In particular convection is paired against
$\varphi\in W^{1,4/3}_{0,\sigma}$ by
$\int_D(u\otimes u):\nabla\varphi$.

\begin{theorem}\label{thm:main}
Let $D\subset\R^3$ be any bounded smooth domain and $\nu>0$.
Let $b_1,\ldots,b_K\in C^\infty(\overline D;\R^3)$ be
divergence-free, tangent to $\partial D$, and satisfy the displayed
positivity condition for some $\delta>0$. On any filtered probability
space satisfying the usual conditions, with independent standard
Brownian motions $W^1,\ldots,W^K$, every deterministic smooth
solenoidal datum $h$ with zero vector trace admits an almost surely
positive stopping time $\tau$ and an adapted solution of
\eqref{eq:NS}, $u(0)=h$, with paths in
\begin{equation}\label{eq:main-class}
 C([0,\tau];B^{1/2}_{4,4}(D;\R^3))
 \cap L^4(0,\tau;W^{1,4}_0(D;\R^3)).
\end{equation}
Any two adapted solutions in this entire class with the same data
and Brownian motions agree up to their common lifetime.

Fixing the domain, coefficients, and Brownian coupling, let smooth
admissible $h_n\to h$ in $\calE$. The corresponding solutions can be
chosen with a positive base time $\tau$ and positive pairwise common
stops $\sigma_n\leq\tau$ such that, for each $\epsilon>0$,
\begin{align}\label{eq:main-dependence}
 \PP\!\left\{\sup_{t\leq\sigma_n}\norm{u_n(t)-u(t)}_{\calE}
 +\left(\int_0^{\sigma_n}\norm{u_n-u}_{V_4}^4\,dt\right)^{1/4}
 >\epsilon\right\}&\longrightarrow0,\\
 \PP\{\sigma_n<\tau\}&\longrightarrow0.\notag
\end{align}
\end{theorem}

Thus perturbations reach the same base interval with probability
tending to one, and converge in the full asserted path topology on
their common intervals. The theorem imposes no amplitude smallness
beyond stochastic parabolicity and no planar-span restriction on
the interior noise. It is local in time; it makes no assertion of
global three-dimensional regularity. The proof uses smooth finite
noise families and does not establish the broader rough or
infinite-noise formulations of the no-slip question.

The main work is linear. In a flat channel, constant tangential
transport is removed by random translation, leaving an anisotropic
Stokes operator. Its pressure and heat roots can coincide even at
a positive resolvent parameter. Divided differences keep the
resolvent regular there. The matrix imposing the two wall conditions
has a uniformly bounded inverse, also after weighting the normal
component. These two bounds preserve the smaller normal trace of the
forcing and the factor $k/a$ in the slow tangential response, where
$k$ and $a$ are the pressure and heat frequency scales. They give the
required operator-valued multiplier bounds. Interior points instead use
a full three-dimensional torus,
so their frozen vectors may have arbitrary rank.

To pass to a curved wall, we localize the Stokes energy form between
the physical and flat domains. This preserves both the no-slip
positive space and its weak dual. Comparing it with Piola transport
requires control of the nonorthogonal pressure term. Negative-order
pressure estimates and the perturbation theorem of Kalton,
Kunstmann, and Weis \cite{KKW} give a quantitative square-root
comparison. Its smallness survives summation over the cover; compact
errors reduce to the physical $L^2$ energy estimate. A finite
continuation in the noise amplitude then constructs the linear
solution operator throughout the parabolicity range. An adapted
state cutoff gives the nonlinear solution, while a separately
stopped energy argument proves uniqueness without moment assumptions
on the competing paths.

Section~\ref{sec:framework} fixes the weak formulation and states the
linear estimate. Section~\ref{sec:models} proves the flat estimates;
Sections~\ref{sec:localization} and \ref{sec:physical-linear} localize
and construct the physical linear solution. Section~\ref{sec:nonlinear}
proves the nonlinear assertions. Throughout, Fourier modes are
$e^{i\xi\cdot x}$, positive Stokes operators have sign $-\PP\Delta$,
and complex square roots have positive real part. Constants may
depend on the fixed domain and coefficients and need not stay
bounded as $\delta\downarrow0$.

\section{Weak Stokes formulation and the linear estimate}\label{sec:framework}

Let $D\subset\R^3$ be bounded with smooth boundary. We use the ordinary
Helmholtz projection $\PP_D$ and set, for $1<q<\infty$,
\[
 H_q=L^q_\sigma(D),\qquad V_q=W^{1,q}_{0,\sigma}(D),
 \qquad Y_q=(V_{q'})^*,\qquad q'=q/(q-1).
\]
Here $H_q$ consists of solenoidal fields with zero weak normal trace;
$V_q$ has zero vector trace. Duals use the integral pairing. Let
$A_D=-\PP_D\Delta$ be the positive strong Dirichlet Stokes operator,
with domain $W^{2,q}(D;\R^3)\cap V_q$. Its weak realization, denoted
$A_{D,w}$, is the form
\begin{equation}\label{eq:weak-stokes}
 \ip{A_{D,w}u}{v}=\int_D\nabla u:\nabla v,
 \qquad u\in V_q,\quad v\in V_{q'}.
\end{equation}
For the equation in this paper the trace space is
\begin{equation}\label{eq:trace-space}
 \calE=(Y_4,V_4)_{3/4,4}
 =\{h\in B^{1/2}_{4,4}(D;\R^3):\divg h=0,\ \Tr h=0\}.
\end{equation}
The trace on the right is the Besov trace. We write $[X_0,X_1]_s$
for complex interpolation and $(X_0,X_1)_{s,p}$ for real interpolation.
All spatial norms are inhomogeneous. The weak Stokes identifications in
\cite[Section 5, Proposition 5.1]{PW} give \eqref{eq:trace-space} and
$[Y_q,V_q]_{1/2}=H_q$. They follow from the isotropic Stokes
square-root scale; they will not be assumed for an anisotropic operator.
The strong isotropic operator has a bounded $H^\infty$ calculus on
$H_q$ by \cite[Theorem 2.1]{KWErratum}, whose $C^{2+\mu}$ boundary hypothesis, $0<\mu<1$, is satisfied by smooth $D$. The negative-half identification
\[
 J_q=(1+A_D)^{-1/2}:Y_q\longrightarrow H_q
\]
is an isomorphism, with $J_q(V_q)=\Dom(A_D)$ and
$J_qA_{D,w}J_q^{-1}=A_D$. Thus the weak operator is similar to the
strong one. We also use the resulting fractional-domain and
interpolation identities for $1+A_D$.

Fix a filtered probability space satisfying the usual conditions and
independent real Brownian motions $W^1,\ldots,W^K$. For predictable
processes write
\[
 \calM_T=C([0,T];\calE)\cap L^4(0,T;V_4),\qquad
 \calX_T=L^4(\Omega;\calM_T),
\]
with sum norms. The noise norm $\gamma(\R^K,H_4)$ is equivalent to
$\norm{(\sum_k|g_k|^2)^{1/2}}_{L^4}$, with constants independent of
finite-dimensional Hilbert indices. No boundary trace is imposed on $g$.
Stochastic integrals below have their usual UMD-space meaning; the
integral estimate used is \cite[Proposition 2.4]{AV}.

Let $b_k\in C^\infty(\overline D;\R^3)$ be solenoidal and tangent to
$\partial D$, and suppose
\begin{equation}\label{eq:parabolicity}
 Q(x):=\nu I-\tfrac12\sum_{k=1}^K b_k(x)\otimes b_k(x)
 \ \geq\delta I>0,\qquad \nu>0.
\end{equation}
Set $B_ku=\PP_D[(b_k\cdot\nabla)u]$. This is bounded
$V_4\to H_4$. The main analytic estimate is the following.

\begin{theorem}\label{thm:linear}
For each fixed $T_0<\infty$, every $0<T\leq T_0$, every
$\mathcal F_0$-measurable $h\in L^4(\Omega;\calE)$, and predictable
\[
 f\in L^4(\Omega\times(0,T);Y_4),\qquad
 g\in L^4(\Omega\times(0,T);\gamma(\R^K,H_4)),
\]
there is a unique predictable $u\in\calX_T$ solving
\begin{equation}\label{eq:linear}
 du+\nu A_{D,w}u\,dt=f\,dt+\sum_k(B_ku+g_k)\,dW^k,
 \qquad u(0)=h,
\end{equation}
as an integral identity in $Y_4$. It satisfies
\begin{equation}\label{eq:linear-bound}
 \norm{u}_{\calX_T}\leq C\left(
 \norm{h}_{L^4(\Omega;\calE)}+
 \norm{f}_{L^4(\Omega\times(0,T);Y_4)}+
 \norm{g}_{L^4(\Omega\times(0,T);\gamma(\R^K,H_4))}\right).
\end{equation}
The constant depends on $D,\nu,\delta$, the fixed finite smooth family
$b$, and $T_0$, and is uniform in $T\leq T_0$. The solution map is
causal and the same bound holds on common stopping intervals.
\end{theorem}

The additive stochastic theorem that we use is
\cite[Theorem 3.14]{AV}: if $X_0$ is isomorphic to a closed subspace
of $L^q$, $q\geq2$, and a shift of $A$ has bounded $H^\infty$
calculus of angle less than $\pi/2$, then $(A,0)$ has stochastic
maximal $L^p$ regularity for $p>2$ and the permitted weights. Here
$p=4$, the weight is zero, $X_0=Y_4$, $\Dom(A)=V_4$, and the
noise space is $[Y_4,V_4]_{1/2}=H_4$. The trace estimate follows from
\cite[Proposition 2.1, Definitions 3.7--3.8]{AV}, using temporal
smoothness $3/8\in(1/4,1/2)$. More explicitly, the estimate supplies
\[
 L^4(0,T_0;V_4)\cap
 H^{3/8,4}(0,T_0;[Y_4,V_4]_{5/8})
 \hookrightarrow C\bigl([0,T_0];([Y_4,V_4]_{5/8},V_4)_{1/3,4}\bigr).
\]
The trace space on the right is $(Y_4,V_4)_{3/4,4}=\calE$, since
$5/8+(1/3)(1-5/8)=3/4$. This gives the $C\calE$ norm as well
as the $L^4V_4$ norm. We always obtain these estimates on a fixed
interval $[0,T_0]$ and then restrict them, so their constants do not
come from a short-interval trace embedding. This is an additive theorem;
the non-small transport in \eqref{eq:linear} requires the arguments below.

\section{Flat transport models}\label{sec:models}

At a boundary point a tangent vector field freezes to a horizontal
constant vector. The corresponding model is a fixed periodic channel.
Interior points require a three-dimensional torus, since the frozen
vectors there may span all of $\R^3$. We first prove the channel
estimate, retaining its pressure equation.

\subsection{Channel resolvent and the colliding roots}\label{sec:channel}

Let $O=(\R^2/\Gamma)\times(0,h)$, where $\Gamma$ is a fixed
full-rank lattice. Write $H_q(O)=L^q_\sigma(O)$,
$V_q(O)=W^{1,q}_{0,\sigma}(O)$, and $Y_q(O)=(V_{q'}(O))^*$,
with periodic horizontal directions and zero vector trace at both walls
in $V_q(O)$. Its horizontal dual frequencies satisfy
$\xi\cdot\gamma\in2\pi\mathbb Z$. Let
\[
 Q=\diag(Q_\tau,\nu),\qquad \delta I\leq Q_\tau\leq\nu I,
 \qquad A_Q=-\PP_O\divg(Q\nabla).
\]
The projection and the pressure gradient are the ordinary ones. One
could normalize diffusion by $x=Ly$, $L=\diag(Q_\tau^{1/2},\sqrt\nu)$,
while leaving velocity components unchanged. Then both gradient and
divergence contain $L^{-1}\nabla_y$, and the lattice becomes
$Q_\tau^{-1/2}\Gamma$. Thus this normalization cannot remove the two
horizontal weights while retaining the original projection. We work in
the original coordinates.

For $\xi\ne0$ put
\begin{equation}\label{eq:roots}
 k=|\xi|,\quad d_\xi=\xi^TQ_\tau\xi,\quad
 \alpha=\sqrt{(\lambda+d_\xi)/\nu},\quad \Rea\alpha>0,
 \quad a=|\alpha|.
\end{equation}
On $|\argu\lambda|\leq\pi-\varepsilon$ we have
$a\geq ck$, $\Rea\alpha\geq ca$, and
$|\lambda|+k^2\asymp a^2$. Constants may depend on the fixed
$\varepsilon,\delta,\nu,h,\Gamma$. The nonzero lattice has $k\geq k_0>0$.
For velocity $(v,w)$ and forcing $(f_\tau,f_3)$ the mode equations are
\begin{align}\label{eq:mode}
 (\lambda+d_\xi-\nu\partial_z^2)v+i\xi p&=f_\tau,&
 (\lambda+d_\xi-\nu\partial_z^2)w+p'&=f_3,\\
 i\xi\cdot v+w'&=0,&(v,w)|_{0,h}&=0.\notag
\end{align}
The transverse component is a scalar Dirichlet resolvent. Eliminating
pressure in the remaining components gives
\begin{equation}\label{eq:clamped}
 \nu(\partial_z^2-k^2)(\partial_z^2-\alpha^2)w
 =k^2f_3+i\xi\cdot f_\tau',\qquad w=w'=0\quad(z=0,h),
\end{equation}
with $v_\parallel=i\xi w'/k^2$ and
$p=-((\lambda+d_\xi-\nu\partial_z^2)w'+i\xi\cdot f_\tau)/k^2$.
At $\xi=0$ the equation is instead $w=0$, $p'=f_3$, and the
tangential scalar Dirichlet problem. This mode has a spectral gap.

The roots $\alpha$ and $k$ collide at $\lambda=\nu k^2-d_\xi$,
which can be positive. Put $\Delta=\alpha^2-k^2$ and use
\[
 \phi_3(z)=\frac{\cosh(\alpha z)-\cosh(kz)}{\Delta},\qquad
 \phi_4(z)=\frac{\sinh(\alpha z)/\alpha-\sinh(kz)/k}{\Delta}.
\]
The removable values are understood analytically. After the two bottom
conditions the remaining boundary matrix is
\[
 M=\begin{pmatrix}\phi_3(h)&\phi_4(h)\\
 \phi_3'(h)&\phi_3(h)\end{pmatrix},\quad
 \det M=\frac{(\alpha/k+k/\alpha)\sinh(\alpha h)\sinh(kh)
 -2(\cosh(\alpha h)\cosh(kh)-1)}{\Delta^2}.
\]
At collision,
\begin{equation}\label{eq:confluent}
 \det M=\frac{\sinh^2(kh)-(kh)^2}{4k^4}>0.
\end{equation}
Indeed the three entries are respectively
$h\sinh(kh)/(2k)$,
$(kh\cosh(kh)-\sinh(kh))/(2k^3)$, and
$(\sinh(kh)+kh\cosh(kh))/(2k)$. The limiting determinant as
$k\to0$ is $h^4/12$; the actual zero mode was treated separately.
For all other parameters, the homogeneous energy identity
\[
 \lambda\norm{u}_2^2+d_\xi\norm{u}_2^2+\nu\norm{u'}_2^2=0
\]
excludes a nonzero solution in the stated sector. Multiplication by
$e^{-i\argu(\lambda)/2}$ also gives the coercive forced estimate.

For the kernel estimates use the divided difference
\[
 D_\alpha(r)=\frac{e^{-\alpha r}-e^{-kr}}{\alpha-k},\qquad
 D_k(r)=-r e^{-kr},\qquad \Pi=\xi\xi^T/k^2.
\]
The decaying half-channel Poisson matrix is
\begin{equation}\label{eq:poisson}
 H_+(r)=\begin{pmatrix}
 e^{-\alpha r}I_2+kD_\alpha(r)\Pi&(i\xi/k)\alpha D_\alpha(r)\\
 i\xi^TD_\alpha(r)&e^{-kr}-kD_\alpha(r)
 \end{pmatrix}.
\end{equation}
It equals $I_3$ at zero and satisfies \eqref{eq:mode} homogeneously.
Let $J=\diag(1,1,-1)$ and define
\[
 R_+=H_+(h),\quad R_-=JH_+(h)J,\quad
 L=\begin{pmatrix}I&R_-\\R_+&I\end{pmatrix},\quad B=L^{-1}.
\]
Thus prescribed boundary values are obtained from
$u(z)=H_+(z)c+JH_+(h-z)Je$ by solving $L(c,e)^T=(g_0,g_h)^T$.

Both $B$ and $\diag(W,W)B\diag(W,W)^{-1}$ are uniformly bounded,
where $W=\diag(I_2,a/k)$. To see this, for comparable $a,k$ use
coercivity and compactness. For $k\to\infty$ both versions of $R_\pm$
tend to zero. For bounded $k$ and $a\to\infty$, putting $\vartheta=a/\alpha$,
\[
 R_+\to e^{-kh}\begin{pmatrix}0&-i\widehat\xi\\0&1\end{pmatrix},
 \qquad WR_+W^{-1}\to e^{-kh}
 \begin{pmatrix}0&0\\-i\vartheta\widehat\xi^T&1\end{pmatrix}.
\]
The reflected limits use $J$. The only possible nonzero eigenvalue
of their product is $e^{-2kh}<1$, uniformly away from one for
$k\geq k_0$. On the remaining compact parameter set, the homogeneous
boundary-value problem has no nonzero solution, so its boundary
matrix is invertible. Consequently every boundary block of $B$ has size
\begin{equation}\label{eq:B-block}
 B_{ee'}=\begin{pmatrix}O(1)&O(1)\\O(k/a)&O(1)\end{pmatrix}.
\end{equation}
The lower-left bound uses the weighted inverse, the upper-right bound
the unweighted inverse. A regular normalized determinant, useful also
at collision, is
\[
 (1-e^{-(\alpha+k)h})^2-(\alpha+k)^2
 \left(\frac{e^{-kh}-e^{-\alpha h}}{\alpha-k}\right)^2.
\]
Its collision value is $4e^{-2kh}(\sinh^2(kh)-(kh)^2)>0$.

\subsection{Kernel derivatives and \texorpdfstring{$R$}{R}-bounds}\label{sec:kernels}

Here all horizontal derivatives keep $\lambda$ fixed. For $c=Q_\tau\xi$,
\begin{equation}\label{eq:root-jets}
 k_i=\xi_i/k,\quad k_{ij}=\delta_{ij}/k-\xi_i\xi_j/k^3,\quad
 \alpha_i=\frac{c_i}{\nu\alpha},\quad
 \alpha_{ij}=\frac{Q_{ij}}{\nu\alpha}-\frac{c_ic_j}{\nu^2\alpha^3}.
\end{equation}
In particular $k^m\partial^\beta\alpha=O(k^2/a)$ for
$m=|\beta|=1,2$. If $E_t(r)=e^{-tr}$, then
\begin{align}\label{eq:exponential-jets}
 \partial_i(t^mE_t)&=E_t(mt^{m-1}-rt^m)t_i,\\
 \partial_{ij}(t^mE_t)&=E_t\bigl[
 (m(m-1)t^{m-2}-2mr t^{m-1}+r^2t^m)t_it_j
 +(mt^{m-1}-rt^m)t_{ij}\bigr].\notag
\end{align}
This also holds for $m=-1$. The angular derivatives are, with $N=\xi\xi^T$,
\begin{align}\label{eq:angular-jets}
 \Pi_i&=(e_i\xi^T+\xi e_i^T)/k^2-2\xi_iN/k^4,\\
 \Pi_{ij}&=(e_ie_j^T+e_je_i^T)/k^2
 -2\{\xi_j(e_i\xi^T+\xi e_i^T)+\xi_i(e_j\xi^T+\xi e_j^T)
 +\delta_{ij}N\}/k^4+8\xi_i\xi_jN/k^6.\notag
\end{align}
For comparable roots put $t_s=(1-s)k+s\alpha$. Then
\begin{align}\label{eq:divided-jets}
 D_\alpha&=-r\int_0^1e^{-t_sr}\,ds,&
 (D_\alpha)_i&=r^2\int_0^1e^{-t_sr}(t_s)_i\,ds,\\
 (D_\alpha)_{ij}&=r^2\int_0^1e^{-t_sr}(t_s)_{ij}\,ds
 -r^3\int_0^1e^{-t_sr}(t_s)_i(t_s)_j\,ds.\notag
\end{align}
For $a\geq2k$ the denominator $|\alpha-k|\geq a/2$ is safe:
if $A=e^{-\alpha r}-e^{-kr}$, $C=\alpha-k$, then
$D_i=(A_i-DC_i)/C$ and
$D_{ij}=(A_{ij}-DC_{ij}-D_iC_j-D_jC_i)/C$.
These formulas, with ordinary product differentiation, give all pure
and mixed second derivatives in \eqref{eq:poisson}.

Let primes now denote distributional derivatives in $z$, and put
\[
 K_\alpha(r)=e^{-\alpha|r|}/(2\alpha),\qquad
 G(r)=\frac{e^{-k|r|}/(2k)-e^{-\alpha|r|}/(2\alpha)}{\Delta}.
\]
The full-line particular velocity kernel is
\begin{equation}\label{eq:full-line}
 R^0(r)=\nu^{-1}\begin{pmatrix}
 K_\alpha(I_2-\Pi)-G''\Pi&i\xi G'\\
 i\xi^TG'&k^2G
 \end{pmatrix}.
\end{equation}
Its crucial distributional identities are
\begin{equation}\label{eq:delta-identity}
 G''''=\delta_0+\frac{k^3e^{-k|r|}-\alpha^3e^{-\alpha|r|}}{2\Delta},
 \qquad K_\alpha''=\alpha^2K_\alpha-\delta_0.
\end{equation}
Thus $\partial_z^2R^0$ contains exactly
$-\nu^{-1}\diag(I_2,0)\delta_0$. The derivatives of its two angular
projections cancel. Away from zero,
$G^{(j)}$ is the signed version of
$(-1)^j(k^{j-1}e^{-k|r|}-\alpha^{j-1}e^{-\alpha|r|})/(2\Delta)$,
with the parity sign on crossing zero.

For $a\geq2k$ use $|\Delta|\geq3a^2/4$ and
\eqref{eq:root-jets}--\eqref{eq:exponential-jets}. Derivatives of
$\Delta^{-1}$ are
$-\Delta_i/\Delta^2$ and
$2\Delta_i\Delta_j/\Delta^3-\Delta_{ij}/\Delta^2$, where
$\Delta_i=2c_i/\nu-2\xi_i$ and
$\Delta_{ij}=2Q_{ij}/\nu-2\delta_{ij}$.
For comparable roots use the nonsingular expression
\[
 G(r)=\frac1{2(\alpha+k)}\int_0^1
 e^{-t_s|r|}\left(\frac{|r|}{t_s}+\frac1{t_s^2}\right)ds.
\]
Since $\Rea t_s\geq ck$, its derivatives of vertical order $j\leq3$
are bounded by $Ck^{j-3}(1+k|r|)^M e^{-ck|r|}$; order four uses
\eqref{eq:delta-identity}. In both regimes the kernels of
\begin{equation}\label{eq:four-families}
 k^{|\beta|}\partial_\xi^\beta
 \{\lambda R^0,\ k^2R^0,\ k\partial_zR^0,\ \partial_z^2R^0\},
 \qquad |\beta|\leq2,
\end{equation}
are bounded by $Cke^{-ck|r|}+Cae^{-ca|r|}$, together with the
identity term just described. Polynomial factors are absorbed by reducing $c$.

For the boundary correction set
\[
 P(z)=[H_+(z),JH_+(h-z)J],\qquad
 T(s)=\operatorname{col}(R^0(-s),R^0(h-s)).
\]
The complete channel resolvent is
\begin{equation}\label{eq:channel-resolvent}
 R(\lambda,\xi;z,s)=R^0(z-s)-P(z)BT(s).
\end{equation}
Put $\mu=|\lambda|+k^2$. At either endpoint, with $d_e(s)$ its distance,
the tangential and normal rows of the trace kernel obey
\begin{align}\label{eq:trace-rows}
 |k^{|\beta|}\partial^\beta T_\tau(s)|
 &\leq C\mu^{-1}(ae^{-ca d_e(s)}+ke^{-ckd_e(s)}),\\
 |k^{|\beta|}\partial^\beta T_3(s)|
 &\leq C\mu^{-1}k(e^{-ca d_e(s)}+e^{-ckd_e(s)}),\quad |\beta|\leq2.
 \notag
\end{align}
These follow separately from $k^2G,kG',G''$ in \eqref{eq:full-line}.
For $a\geq2k$, the Poisson blocks have sizes
\begin{equation}\label{eq:poisson-sizes}
 \begin{aligned}
 H_{\tau\tau}&=O(E_a+(k/a)E_k),& H_{\tau3}&=O(E_a+E_k),\\
 H_{3\tau}&=O((k/a)(E_a+E_k)),& H_{33}&=O(E_k+(k/a)E_a),
 \end{aligned}
\end{equation}
where $E_b=e^{-cb r}$. They retain these sizes under normalized
horizontal derivatives. In particular the slow response to a tangential
input has the factor $k/a$.

For completeness, inverse differentiation gives
\[
 B_i=-BL_iB,\qquad
 B_{ij}=BL_jBL_iB+BL_iBL_jB-BL_{ij}B.
\]
Differentiating $WLW^{-1}$ and using
$k\partial_iW\,W^{-1}$ and its second analogues, which are bounded
by \eqref{eq:root-jets}, proves \eqref{eq:B-block} for all the
normalized first and second horizontal derivatives of $B$. The second
derivative of the reflected
product has all nine terms
\begin{align}\label{eq:nine-terms}
 &(PBT)_{ij}=P_{ij}BT+P_iB_jT+P_iBT_j+P_jB_iT+PB_{ij}T\notag\\
 &\hspace{35mm}{}+PB_iT_j+P_jBT_i+PB_jT_i+PBT_{ij}.
\end{align}
For vertical order $j_0=0,1,2$ replace $P$ by $\partial_z^{j_0}P$;
multiply by $\lambda$ or by $k^{2-j_0}$ as appropriate in
\eqref{eq:four-families}. Fast parts gain $a^{j_0}$, slow parts
$k^{j_0}$. The fast factor $k^{2-j_0}a^{j_0}\leq Ca^2$ cancels
$\mu^{-1}$. Whenever an $a$-sized tangential trace enters a slow
output, \eqref{eq:poisson-sizes} and \eqref{eq:B-block} supply $k/a$.
For example, a fast tangential trace entering a slow tangential output
in $\lambda PBT$ is bounded by
\[
 C|\lambda|\frac{k}{a}\frac{a}{\mu}
 e^{-ckd_e(z)-ca d_{e'}(s)}
 \leq Ck e^{-c'k(d_e(z)+d_{e'}(s))}
 \leq\frac{C}{d_e(z)+d_{e'}(s)}.
\]
Here $|\lambda|/\mu\leq1$ and $a\geq ck$. For the second spatial
families the slow output contributes $k^2$, giving
$k^3/\mu\leq Ck$ instead. Fast--fast terms are bounded by
$Ca e^{-ca(d_e(z)+d_{e'}(s))}$; a fast output with a slow input
has size at most $Ck$ with decay at least at scale $k$ on both sides.
The block bounds for $B$ preserve these estimates, and the smaller
normal row in \eqref{eq:trace-rows} handles the remaining columns.
The same reasoning applies to the normalized derivatives just listed.
Every resulting boundary kernel is bounded by
\begin{equation}\label{eq:hardy-bound}
 C\sum_{e,e'\in\{0,h\}}
 \frac1{d_e(z)+d_{e'}(s)}.
\end{equation}
For comparable roots $a\asymp k$, \eqref{eq:divided-jets} gives
polynomials in $kd_e(z)$ and $kd_{e'}(s)$ times the same
exponentials; absorbing the polynomials by a smaller exponential
constant gives the same estimate.
This accounts for pure and mixed jets, the reflected product, and the
distributional identity; individual mode energy bounds alone would not suffice.

Each positive kernel in \eqref{eq:hardy-bound} is bounded on $L^q(0,h)$:
split $1/(z+s)$ at $s=z$ into the two Hardy operators, and reflect for
the other endpoints. Minkowski's inequality in $\ell^2$ therefore
gives the same bound for diagonal families. The full-line exponential
kernels are dominated by the one-dimensional Hardy--Littlewood maximal
operator. Its vector-valued estimate
\cite[Theorem 1]{DK}, with outer exponent $q$ and sequence exponent two,
gives the corresponding $\ell^2$ bounds. These arguments work on
$L^q(0,h;\ell^2_N)$ uniformly in $N$ and with another Hilbert index.
Recall that a family $\mathcal T\subset\mathcal L(X)$ is
$R$-bounded if, for independent symmetric signs $r_\ell$,
\[
 \left(\EE_r\left\Vert\sum_\ell r_\ell T_\ell x_\ell\right\Vert_X^2\right)^{1/2}
 \leq C\left(\EE_r\left\Vert\sum_\ell r_\ell x_\ell\right\Vert_X^2\right)^{1/2}
 \quad(T_\ell\in\mathcal T).
\]
On $L^q$, Khintchine's inequality identifies these randomized norms
with the corresponding $\ell^2$ square-function norms. The preceding
bounds therefore show that the families in \eqref{eq:four-families},
now for $R$, are $R$-bounded, including arbitrary finite collections
of $\lambda$'s.

\subsection{Domains, multipliers, and stochastic estimates}\label{sec:channel-consumer}

The multiplier theorem \cite[Corollary 4.4]{GW} applies to UMD spaces
of Fourier type $r>1$ when all normalized total derivatives through
$\lfloor n/r\rfloor+1$ are $R$-bounded. The spaces
$L^q(0,h;\C^3\otimes\ell^2_N)$ have Fourier type
$\min(q,q')$ and UMD constants uniform in $N$, by the Rademacher
embedding into $L^q$; see \cite[Remarks 2.3, 2.6, 2.10]{GW}.
In horizontal dimension two, second derivatives suffice for every
fixed finite $q>1$. Apply the theorem to the diagonal symbol with
entries $\lambda_jR(\lambda_j,\xi)$. The preceding estimates have the
additional Hilbert-index uniformity required to do this. We obtain
the $R$-resolvent bound on $L^q(\R^2\times(0,h))$, and the spatial
domain bounds from the other three families.

Here is the periodic restriction needed for this application. First
take the standard lattice; the fixed linear change for general
$\Gamma$ is addressed below. Cut off
the symbol smoothly below $k_0/2$, with value one at every nonzero
lattice frequency. For a finite trigonometric polynomial
$F(x)=\sum_j e^{ij\cdot x}x_j$, multiply by $\phi(\epsilon x)$,
where $\widehat\phi$ is smooth and compactly supported. For each
frequency $j$, after $y=\epsilon x$, the Fourier error is
$[m(j+\epsilon\xi)-m(j)]\widehat\phi(\xi)$.
It tends uniformly to zero with its first two derivatives. Two
integrations by parts bound its inverse transform by
$o(1)(1+|y|)^{-2}$, whose $q$th power is integrable in $\R^2$.
Thus the error in the original $L^q$ norm is $o(\epsilon^{-2/q})$.
Periodic averaging transfers the Euclidean multiplier bound to the
samples. The constants are uniform in the Hilbert index; density
finishes the restriction. A fixed lattice change uses the pure as
well as mixed second derivatives already proved. Add the separately
solved scalar zero mode.

At $\lambda=1$ the spatial bounds imply that the constructed resolvent
maps full $L^q$ forcing into
$\Dom_q=W^{2,q}(O;\C^3)\cap V_q(O)$. Its pressure gradient follows
from \eqref{eq:mode}. Conversely $f=(1+A_Q)u$ for $u\in\Dom_q$
has this solution by mode uniqueness. It is the ordinary projected
resolvent on solenoidal forcing, so $A_Q$ is closed, densely defined,
and has exactly this domain. The zero-mode gap and the nonzero-mode
bounds give invertibility at zero. Resolvent sectors give a causal
analytic semigroup. For the time symbol
$M(t)=A_Q(it+A_Q)^{-1}$, both $M$ and
$tM'=-itA_Q(it+A_Q)^{-2}$ are $R$-bounded. The dimension-one case
of \cite[Corollary 4.4]{GW} gives deterministic maximal $L^p$
regularity, $1<p<\infty$. Semigroup causality and zero extension of
forcing give constants uniform in $T\leq T_0$. In particular this
holds at $q=4,4/3,8,8/7$.

We next need the isotropic channel calculus; a bounded Euclidean-domain
Stokes theorem does not directly have this geometry. Reflect tangential
components evenly and normal components oddly into the doubled torus,
apply its Helmholtz projection, and restrict. The result fixes
solenoidal zero-normal fields and annihilates gradients, hence is
$\PP_O$. Fourier multipliers give its $L^q$ bounds and consistent
adjoints, and reflection gives $H^{s,2}$ bounds for $0<s<1/2$.
Let $A_0$ be the isotropic channel Stokes operator. The elliptic
projection
\[
 \widetilde Q=A_0^{-1}\PP_O(-\Delta_D)
\]
is bounded on full $L^q$: its adjoint is
$(-\Delta_D)A_{0,q'}^{-1}\PP_O$, bounded by the strong-domain estimate.
It fixes $H_q$ and maps full-vector $W^{2,q}\cap W^{1,q}_0$ into
$\Dom_q$. Interpolating this retraction gives
$[H_q,\Dom_q]_{1/2}=V_q$. The componentwise Dirichlet Laplacian is
the odd part of the doubled-torus Laplacian and has its scalar
$H^\infty$ calculus.

We use the corrected comparison theorem
\cite[Theorem 1.3]{KWErratum}. Let $(X_0,X_1)$ be a reflexive
B-convex interpolation couple, with compatible surjections
$P_j:X_j\to Y_j$ and compatible bounded right inverses
$\mathsf J_j:Y_j\to X_j$. Let $\mathsf A_j,\mathsf B_j$ be the
consistent sectorial families on these spaces. Suppose that
$\mathsf A_j$ has bounded $H^\infty$ calculus and $\mathsf B_j$ is
almost $R$-sectorial. For $\alpha<0<\beta$, require both
\begin{equation}\label{eq:corrected-images}
 P_0\dot X_{0,\alpha,\mathsf A_0}
   =\dot Y_{0,\alpha,\mathsf B_0},\qquad
 P_1\dot X_{1,\beta,\mathsf A_1}
   =\dot Y_{1,\beta,\mathsf B_1},
\end{equation}
and the compatible extensions
\begin{equation}\label{eq:corrected-right-inverses}
 \mathsf J_0:\dot Y_{0,\alpha,\mathsf B_0}
                 \longrightarrow\dot X_{0,\alpha,\mathsf A_0},\qquad
 \mathsf J_1:\dot Y_{1,\beta,\mathsf B_1}
                 \longrightarrow\dot X_{1,\beta,\mathsf A_1}.
\end{equation}
Then $\mathsf B$ has bounded $H^\infty$ calculus on
$[Y_0,Y_1]_\theta$, $0<\theta<1$. The dots denote the homogeneous
fractional scales of the specified operators. The estimates
\eqref{eq:corrected-right-inverses} are additional hypotheses; the
surjections in \eqref{eq:corrected-images} do not imply them.

Here $A_{0,p}=-\PP_O\Delta$ denotes the unit-viscosity reference
Stokes operator, and $L_p=-\Delta_{\mathrm{Dir},p}$ the componentwise
Dirichlet Laplacian. The physical viscosity remains $\nu$ in
\eqref{eq:linear} and in $A_Q$. With $i_p:H_p\hookrightarrow L^p$,
set
\begin{equation}\label{eq:channel-right-inverse}
 \mathcal J_p=L_pi_pA_{0,p}^{-1},\qquad
 \widetilde Q_p=A_{0,p}^{-1}\PP_OL_p.
\end{equation}
The new symbol $\mathcal J_p$ denotes this elliptic right inverse,
not the inverse-root isomorphism $J_q$ of Section~\ref{sec:framework}.
The strong-domain estimate already proved above makes
$\mathcal J_p:H_p\to L^p$ bounded and gives
$\PP_O\mathcal J_p=I$. Its fractional estimate requires one more
regularity argument.

\begin{lemma}\label{lem:channel-extra-derivative}
On the fixed channel, let $f\in H^1(O;\C^3)$ and
$u=A_{0,2}^{-1}\PP_Of$. Choose the periodic pressure $p$ with mean
zero in $-\Delta u+\nabla p=f$. Then
\begin{equation}\label{eq:channel-h3}
 \norm{u}_{H^3}+\norm{p}_{H^2}\leq C\norm{f}_{H^1}.
\end{equation}
The constant depends on the fixed lattice and height.
\end{lemma}
\begin{proof}
The strong-domain bound and the Helmholtz decomposition first give
\begin{equation}\label{eq:channel-h2-pressure}
 \norm{u}_{H^2}+\norm{p}_{H^1}\leq C\norm{f}_{L^2}.
\end{equation}
For a horizontal direction $e$, take the difference quotient
$\delta_t^e$. Translation preserves the channel, the zero wall values,
and the pressure normalization. Thus \eqref{eq:channel-h2-pressure}
applied to the translated difference gives
\[
 \norm{\delta_t^e u}_{H^2}+\norm{\delta_t^e p}_{H^1}
 \leq C\norm{\delta_t^e f}_{L^2}
 \leq C\norm{\partial_e f}_{L^2}.
\]
Weak compactness and the distributional difference-quotient limit yield
$\partial_eu\in H^2$ and $\partial_ep\in H^1$. This controls every
third velocity derivative containing a horizontal derivative, and every
second pressure derivative containing one.

Write $u=(u_h,u_3)$ and use $z$ for the normal coordinate.
Incompressibility and the normal momentum equation give, successively,
\begin{equation}\label{eq:channel-normal-reconstruction}
 \partial_z^3u_3=-\divg_h\partial_z^2u_h,\qquad
 \partial_z^2p=\partial_zf_3+\Delta\partial_zu_3.
\end{equation}
Their right-hand sides are controlled in $L^2$ by the preceding
horizontal estimates and $\norm{f}_{H^1}$. The horizontal momentum
equations then give
\[
 \partial_z^3u_h=-\partial_zf_h-\Delta_h\partial_zu_h
                         +\nabla_h\partial_zp.
\]
These are the remaining third velocity derivatives. Together they
prove \eqref{eq:channel-h3}. The identities are distributional;
no trace of $f$, $\Delta u$, or $p$ has been prescribed at the wall.
\end{proof}

\begin{lemma}\label{lem:channel-compatible-right-inverses}
The operators $\mathcal J_p$ are compatible base-space right inverses
of $\PP_O$. They have consistent bounded extensions
\begin{align}
 \mathcal J_2:&\ \Dom(A_{0,2}^{s})\longrightarrow\Dom(L_2^{s}),
                     &&0<s<\tfrac14,\label{eq:channel-positive-J}\\
 \mathcal J_p:&\ Z_{A,p}\longrightarrow Z_{\Delta,p},
                     &&1<p<\infty,\label{eq:channel-negative-J}
\end{align}
where
\[
 Z_{\Delta,p}=\bigl(\Dom(L_{p'})\bigr)^*,\qquad
 Z_{A,p}=\bigl(\Dom(A_{0,p'})\bigr)^*.
\]
Moreover $\PP_O\mathcal J_p=I$ on these extended spaces.
\end{lemma}
\begin{proof}
The closed coercive form $\int_O\nabla u:\overline{\nabla v}$ on
$V_2$ has positive selfadjoint operator $A_{0,2}$, as follows by
identifying its resolvent with the strong solution already constructed.
The Hilbert form theorem gives
$\Dom(A_{0,2}^{1/2})=V_2$; see \cite[Theorem 2.13 and (3.51)]{Teschl}.
This Hilbert assertion precedes the $L^p$ calculus under consideration.
The base estimate bounds $\mathcal J_2:H_2\to L^2$, and
Lemma~\ref{lem:channel-extra-derivative} bounds
$\mathcal J_2:V_2\to H^1$, since $L_2u\in H^1$ when $u\in H^3$.
Interpolating the same operator at parameter $2s$ gives
\begin{equation}\label{eq:channel-J-interpolation}
 \mathcal J_2:
 [H_2,\Dom(A_{0,2}^{1/2})]_{2s}
   =\Dom(A_{0,2}^{s})
 \longrightarrow [L^2,H^1]_{2s}=H^{2s},\qquad 0<s<\tfrac12.
\end{equation}
The fractional-domain identity here follows from the Hilbert spectral
calculus and \cite[Proposition 2.2]{KKW}, not from the desired
Banach-space calculus.

For $0<s<1/4$, one also has
\begin{equation}\label{eq:channel-scalar-subtrace}
 \Dom(L_2^s)=H^{2s}(O;\C^3).
\end{equation}
Indeed, use finitely many horizontal periodic coordinate cutoffs and
zero extension of the corresponding flat strip pieces. The indicator
bound \cite[Proposition 2.4]{BG} applies with $p=2$ and
$0<2s<1/2$. Odd reflection at the walls therefore extends these
pieces boundedly into the doubled periodic space; restriction is its
bounded inverse on the odd subspace. Fourier series horizontally and
sine series vertically identify its weighted norm with
$\norm{L_2^s\cdot}_{L^2}$. Ordinary Sobolev interpolation in
\eqref{eq:channel-J-interpolation} follows similarly using a common
even extension at orders zero and one and Fourier-weight interpolation.
Equations~\eqref{eq:channel-J-interpolation}--\eqref{eq:channel-scalar-subtrace}
prove \eqref{eq:channel-positive-J}. The strict subtrace range avoids
an extra boundary condition on $\mathcal J_2f$.

For the negative endpoint, $\widetilde Q_{p'}$ is bounded from
$\Dom(L_{p'})$ into $\Dom(A_{0,p'})$. Define
\[
 \langle\mathcal J_p\ell,v\rangle
   =\langle\ell,\widetilde Q_{p'}v\rangle,
 \qquad v\in\Dom(L_{p'}),\quad\ell\in Z_{A,p}.
\]
This proves \eqref{eq:channel-negative-J} and agrees with the base
operator by adjointness. The projection extends by restricting
functionals to the continuously included domain
$\Dom(A_{0,p'})\subset\Dom(L_{p'})$. Since
$\widetilde Q_{p'}v=v$ on the solenoidal strong domain,
$\PP_O\mathcal J_p\ell=\ell$. The same identity on base spaces follows
from $-\PP_O\Delta A_{0,p}^{-1}=I$. Consistency follows from the common
elliptic solutions and the distributional formulas; thus these are
extensions of one compatible family.
\end{proof}

We now verify the projection images as well. The retraction
$\widetilde Q_2$ maps the scalar base and strong Dirichlet domains to
the Stokes counterparts. Interpolating these Hilbert endpoint bounds
maps $\Dom(L_2^s)$ into $\Dom(A_{0,2}^s)$. Since $\PP_O$ is bounded
on $H^{2s}$ below the trace threshold, and
$\widetilde Q_2\PP_O=\PP_O$, this proves its bounded positive image.
Lemma~\ref{lem:channel-compatible-right-inverses} makes it onto:
\[
 \PP_O\Dom(L_2^s)=\Dom(A_{0,2}^s),\qquad 0<s<\tfrac14.
\]
The same lemma gives the negative full-step image
$\PP_OZ_{\Delta,p}=Z_{A,p}$, together with its right inverse.

Apply the corrected criterion to inverse operators with
\[
 p_0=2,\qquad p_1=8\ \hbox{or}\ 8/7,\qquad
 \mathsf A_j=L_{p_j}^{-1},\qquad
 \mathsf B_j=A_{0,p_j}^{-1},\qquad
 P_j=\PP_O,\quad\mathsf J_j=\mathcal J_{p_j}.
\]
Take $\alpha=-1/8$ and $\beta=1$. Inversion exchanges the fractional
signs, so \eqref{eq:channel-positive-J} and
\eqref{eq:channel-negative-J} are exactly the two estimates in
\eqref{eq:corrected-right-inverses}. The scalar calculus, compatible
resolvents, and $R$-sectoriality were proved before this application.
At $\theta=2/3$ the resulting spaces are $H_4$ and $H_{4/3}$.
Calculus for the inverse is equivalent to calculus for $A_0$ by
$F(z)\mapsto F(1/z)$. This first establishes calculus existence.
Only then do \cite[Theorem 4.1, Corollary 4.4]{KKW} and the proved
$R$-sectorial angle zero give a calculus angle less than $\pi/2$.

The resulting negative-half scale identifies
$Y_q=(V_{q'})^*$ and gives $[Y_q,V_q]_{1/2}=H_q$.
For general $Q$, the same causal zero-initial solution map has the
consistent endpoint bounds, for $1<p<\infty$,
\[
 \begin{aligned}
 L^p(0,T;H_q)&\longrightarrow
 W^{1,p}(0,T;H_q)\cap L^p(0,T;\Dom_q),\\
 L^p(0,T;Z_q)&\longrightarrow
 W^{1,p}(0,T;Z_q)\cap L^p(0,T;H_q),
 \qquad Z_q=(\Dom_{q'})^*.
 \end{aligned}
\]
The second line is the one-step extrapolation, equivalently obtained
by duality from the adjoint strong estimate. The isotropic retraction
identities give
\[
 [Z_q,H_q]_{1/2}=Y_q,\qquad [H_q,\Dom_q]_{1/2}=V_q.
\]
Interpolate the solution and its time derivative at parameter $1/2$.
This proves deterministic maximal regularity of the weak form with
domain $V_q$ on $Y_q$. Integration by parts identifies the form with
$\int Q\nabla u:\nabla v$. No general-$Q$ root identity is required.
The real trace is $(Y_4,V_4)_{3/4,4}=B^{1/2}_{4,4,0,\sigma}$,
obtained by interpolating isotropic levels on either side of spatial
smoothness $1/2$, both above the trace threshold $1/4$.

For this model set
$\calE_O=(Y_4(O),V_4(O))_{3/4,4}$ and
$\calX_T(O)=L^4(\Omega;C([0,T];\calE_O)\cap L^4(0,T;V_4(O)))$.
The additive stochastic theorem stated in Section~\ref{sec:framework}
applies to the isotropic weak operator: $Y_4$ is similar to the closed
subspace $H_4$ of $L^4$. For general $Q$, solve the zero-initial
isotropic stochastic equation for $z_g$, then correct it with the
zero-initial deterministic $A_Q$ solution forced by
$(A_{0,w}-A_{Q,w})z_g+f$. The form difference is bounded
$V_4\to Y_4$. For nonzero $h\in\calE$, first extend by
$z_h=e^{-tA_{0,w}}h$ and use the same deterministic correction on
$(A_{0,w}-A_{Q,w})z_h$. The isotropic real trace estimate bounds
$z_h$ in $C\calE\cap L^4V_4$. Both corrections have this trace by
deterministic maximal regularity. These maps are causal and preserve
predictability. They prove the full additive channel estimate, including
random $\mathcal F_0$-measurable initial data, on the required spaces.

\subsection{Interior torus and removal of constant transport}\label{sec:torus}

On a fixed full three-dimensional torus put
$H_q=L^q_\sigma$, $V_q=W^{1,q}_\sigma$, $Y_q=(V_{q'})^*$,
and $\calE_O=(Y_4,V_4)_{3/4,4}=B^{1/2}_{4,4,\sigma}$.
These norms include constants. For arbitrary constant vectors
$\beta_1,\ldots,\beta_K\in\R^3$ set
\[
 T_k=\beta_k\cdot\nabla,\qquad
 Q=\nu I-\tfrac12\sum_k\beta_k\otimes\beta_k\geq\delta I,
 \qquad A_Q=-\divg(Q\nabla).
\]
Here $\PP(0)=I$, and $\PP(\xi)=I-\xi\xi^T/|\xi|^2$ otherwise.
Let $J_0=(1-\Delta)^{1/2}$. It is an isomorphism from $H_q$ to
$Y_q$, and $J_0^{-1}$ sends $V_q$ to $W^{2,q}_\sigma$.
For bounded holomorphic $F$, Cauchy's estimates give all normalized
frequency derivatives through order four of
$F(1+\xi^TQ\xi)$. Scalar multipliers, or the scalar case of
\cite[Corollary 4.4]{GW}, give a bounded $H^\infty$ calculus for
$1+A_Q$ on $H_q$, with domain $W^{2,q}_\sigma$. Similarity by
$J_0^{-1}$ gives its weak calculus with domain $V_q$ on $Y_q$.

The three-dimensional restriction can be proved by the same window
argument as above. This time differentiate the compact Fourier error
through order four and apply $(1-\Delta_\xi)^2$. The inverse transform
is $o(1)(1+|y|^2)^{-2}$, integrable to the $q$th power for
$q=4,4/3$ in $\R^3$. Periodic averaging and density finish the proof
for a fixed, possibly rotated lattice. The scalar multiplier is smooth
at zero; the projection's constant mode is kept separately.
The calculus constants are uniform for $\delta I\leq Q\leq\nu I$.
The additive stochastic theorem therefore applies on $Y_4$.
One fixed shift handles the mean and is removed by an exponential
change of unknown on $[0,T_0]$.

On either flat geometry consider
\begin{equation}\label{eq:flat-transport}
 du+A u\,dt=f\,dt+\sum_k(T_ku+g_k)\,dW^k,
 \qquad A=-\nu\PP_O\Delta.
\end{equation}
In the channel the $\beta_k$ are horizontal; on the torus they are
arbitrary. Translations
$U_tv(x)=v(x+\sum_k\beta_kW^k(t))$ are isometries of all the
spaces just defined and commute with $A,T_k$. Since
$A_Q=A+\frac12\sum_kT_k^2$, It\^o's product rule gives
\begin{equation}\label{eq:translation}
 dv+A_Qv\,dt=
 \left(U_t^{-1}f-\sum_k T_kU_t^{-1}g_k\right)dt
 +\sum_kU_t^{-1}g_k\,dW^k,\qquad v=U_t^{-1}u.
\end{equation}
The correction involving $g$ is essential. Each $T_k:H_4\to Y_4$
is bounded by integration by parts. Thus all the transformed data
are in the additive data spaces. Smooth horizontal regularization in
the channel, and Fourier regularization on the torus, commute with
these operators. They prove \eqref{eq:translation} first for smooth
approximants. The $\calX_T$ and stochastic integral bounds then pass
to the weak equation for rough data. This proves the flat transport
estimate, uniqueness, and causality.

On the torus the mean satisfies exactly
\[
 \overline u(t)=\overline h+\int_0^t\overline f(s)\,ds
 +\sum_k\int_0^t\overline g_k(s)\,dW^k(s).
\]
Its finite-dimensional estimate controls its contribution to the full
$W^{1,4}$ norm. There is no restriction on the rank of the family
$\beta$. For a common stopping interval, extend only the inhomogeneous
data by zero to $[0,T_0]$ and restrict the causal solution. The
multiplicative coefficients still define the full-interval equation;
uniqueness identifies the restrictions. Constants are uniform in
$T\leq T_0$ and, for bounded $\beta$, in all covariances between
$\delta I$ and $\nu I$.

\section{Localization on the weak Stokes scale}\label{sec:localization}

A cutoff followed by Helmholtz projection need not preserve zero vector
trace. To localize simultaneously on $V_q$ and $Y_q$ we instead use
the isotropic Stokes energy form. All roots in this section belong
to isotropic operators. For either model domain $O$ we use
$H_q(O),V_q(O),Y_q(O)$ as above, with no wall trace on the full
torus, and write $\calE_O=(Y_4(O),V_4(O))_{3/4,4}$.

\subsection{Energy-form localizers}\label{sec:energy-localizers}

For a physical domain $D$ and a model $O$ write $M_D=1+A_D$,
$M_O=1+A_{0,O}$, and use the same letters for their consistent
extrapolations. On the torus the shift includes constants. Let
$\mathsf E$ be a smooth cutoff and coordinate map on full vector
fields. Suppose $\mathsf E,\mathsf E^*$ are $L^q$ bounded and
preserve full-vector $W^{1,q}_0$ and $W^{2,q}\cap W^{1,q}_0$.
The latter means zero value trace, without a condition on derivative
traces. The actual charts below have these properties for each fixed
radius. Define
\begin{equation}\label{eq:cross-form}
 \ip{Cu}{v}=\int_O\bigl(\nabla(\mathsf E u):\nabla v
 +(\mathsf E u)\cdot v\bigr),\qquad
 T=M_O^{-1/2}CM_D^{-1/2}.
\end{equation}
The form is bounded $V_q(D)\to Y_q(O)$, so $T:H_q(D)\to H_q(O)$.
For $u\in\Dom(A_D)$ its strong representative is
$Cu=\PP_O(1-\Delta_O)\mathsf E u\in H_q(O)$. Since
$M_D^{-1/2}:V_q\to\Dom(A_D)$, this proves $T:V_q(D)\to V_q(O)$.
Similarly
\[
 C^*v=\PP_D\mathsf E^*(1-\Delta_O)v,
 \qquad v\in\Dom(A_{0,O}),
\]
gives $T^*:V_{q'}(O)\to V_{q'}(D)$. Duality extends $T$ consistently
to $Y_q(D)\to Y_q(O)$. We use both $q=4$ and $q=4/3$.

Define the elliptic projection on full vector fields by
\begin{equation}\label{eq:elliptic-projection}
 Q_O=M_O^{-1}\PP_O(1-\Delta_O),\qquad S=Q_O\mathsf E.
\end{equation}
Its full $L^q$ extension follows by duality from the strong Stokes
estimate, as in Section~\ref{sec:channel-consumer}. It fixes
solenoidal fields and preserves zero-value-trace orders one and two.
Equations \eqref{eq:cross-form}--\eqref{eq:elliptic-projection} give
\begin{equation}\label{eq:T-S}
 T=M_O^{1/2}SM_D^{-1/2}.
\end{equation}
We will compare this map with ordinary Piola localization, rather than
assume that the latter acts on the weak dual scale.

\subsection{Geometry, divergence correction, and pressure}\label{sec:geometry}

Choose a regular radius-$r$ cover of $D$ with bounded overlap,
$N\leq Cr^{-3}$, and partition $\sum_j\chi_j=1$. The two model
cells are fixed independently of $r$: a full three-dimensional torus
for interior patches and a periodic channel for boundary patches.
Interior charts are affine rotations. Boundary charts, centered on
$\partial D$, are volume-preserving graph shears
\begin{equation}\label{eq:charts}
 F_j(y)=x_j+O_j(y',y_3+\psi_j(y')),\quad O_j\in SO(3),\quad
 G_j=DF_j,\quad R_j=G_j^{-1},\quad
 \mathsf E_ju=\chi_jR_j(u\circ F_j).
\end{equation}
In the chart formula, $\chi_j$ denotes the pulled-back cutoff
$\chi_j\circ F_j$. Cutting off the graph outside a $Cr$ neighborhood gives
$\norm{\psi_j}_\infty\leq Cr^2$,
$\norm{\nabla\psi_j}_\infty\leq Cr$,
$\norm{D^2\psi_j}_\infty\leq C$, and uniform bi-Lipschitz bounds.
Also $|\partial^m\chi_j|\leq C_mr^{-m}$. Smooth composition and
multiplication show the value-trace and strong-domain properties
used above for both $\mathsf E_j$ and its integral adjoint.
The frozen vectors are $\beta_{j,k}=O_j^Tb_k(x_j)$. They are
horizontal at a boundary center and arbitrary at an interior center.
Orthogonal freezing preserves \eqref{eq:parabolicity}.

We need a negative-order pressure estimate to compare the physical and
flat Laplacians. Put
\begin{equation}\label{eq:buffers}
 \eta=1/8,\qquad d=1/16,\qquad \zeta=3/16,\qquad
 \theta=1-d=15/16.
\end{equation}
These are fixed proof parameters. The value $\eta$ lies strictly
inside both pressure and subtrace windows, $2d=\eta$ matches spatial
and fractional orders, and $\theta=1-d$ yields a lower-order gain
$\rho^{-d}$. The margin $\eta<\zeta<1/4$ permits multiplication
at order $-\eta$ and leaves
$1-\zeta-3/4=1/16>0$ after summing the cover; see
Section~\ref{sec:reconstruction}.
The ordinary Neumann gradient projection in
\cite[Theorem 4.38, Proposition 4.39]{BG} is bounded on $H^{s,q}$
for $-1+1/q<s<1/q$ on smooth bounded domains. Take $s=\pm\eta$
and $q=4,4/3$, strictly inside these intervals. If
$\nabla p_u=(I-\PP_D)\Delta u$, fixing its mean on each component,
the projection bound and the Ne\v cas estimate
\cite[Lemma 2.12]{BG} give
\begin{equation}\label{eq:negative-pressure}
 \norm{p_u}_{H^{1-\eta,q}/\R}
 +\norm{\nabla p_u}_{H^{-\eta,q}}
 \leq C\norm{\Delta u}_{H^{-\eta,q}}.
\end{equation}
The potential estimate extends from smooth fields by density.
The cited result is used in its fixed preprint version.
For $G=D$ or a model $O$, use the isotropic scale
$X_{s,G}=\Dom(M_{G,q}^s)$ for $s\geq0$, and
$X_{-s,G}=(\Dom(M_{G,q'}^s))^*$ for $s>0$, with the integral
pairing extending the $H_q,H_{q'}$ duality. This is equivalently the
completion in the norm $\norm{M_{G,q}^{-s}\,\cdot\,}_{H_q}$.
The isotropic fractional scale then gives
\[
 \Dom(M_D^\theta)=H^{2-\eta,q}\cap V_q,
 \qquad X_{-d,O}=\PP_OH^{-\eta,q},
\]
the latter consistently by duality from
$\Dom(M_{O,q'}^d)=H^{\eta,q'}_\sigma$, below the trace thresholds.
Canonical zero extension at order $-\eta$ is bounded by
\cite[Proposition 2.4]{BG}; its strict window contains our exponents.
In what follows we extend $\partial_{il}u$ and $\nabla p_u$
themselves, rather than differentiate the zero extension of $u$.

Fix one convex Euclidean box in each model coordinate cell, the
boundary box having its bottom on the wall. It contains every small
cutoff of that type. The regularized supported divergence solver of
\cite[formula (3.13), Theorem 3.2, Corollary 3.4, Remark 3.5]{CM}
uses one smoothing kernel on each box and satisfies
\begin{equation}\label{eq:bogovskii}
 \divg\mathcal R_B g=g-\vartheta\int g,\qquad
 \mathcal R_B:H^{s,q}_{\mathrm{supp}}\to
 H^{s+1,q}_{\mathrm{supp}},\quad s=-1/8,0,1.
\end{equation}
Here $\int\vartheta=1$. The same operator acts consistently at all
these exponents and orders. Its supported outputs extend into the
model; at order one they have zero value trace, and at order two
the supported extension also has zero derivative traces. We use the
Euclidean result on these fixed boxes, not a presumed torus version.

For solenoidal $f$, the Piola identity gives
$\divg\mathsf E_jf=\nabla\chi_j\cdot R_j(f\circ F_j)$, of mean
zero. This follows first from the weak zero-normal identity and then
by density on $X_{-d,D}$, using the consistent distributional
extensions of the same Piola map and divergence correction. Consequently,
suppressing $j$,
\begin{align}\label{eq:divergence-comparison}
 (\PP_O-Q_O)\mathsf E f
 &=(\PP_O-Q_O)\mathcal R_B\divg\mathsf E f,\\
 (S-\mathsf E)f&=(Q_O-I)\mathcal R_B\divg\mathsf E f.\notag
\end{align}
On the channel, even tangential and odd normal reflection, followed
by full-torus multipliers, bounds $\PP_O$ on the supported positive
spaces here; zero vector trace makes reflection bounded through
order one. On the torus use its Fourier projection. Thus, for fixed $r$,
$S-\mathsf E:H_q\to W^{1,q}$ and $S-\mathsf E:V_q\to W^{2,q}$.
It is compact at orders zero and one, respectively. For an
$H^{-1/8,q}$ input, the same correction gains one order before
$\PP_O,Q_O$ are applied.

\subsection{Two defect estimates and square roots}\label{sec:roots}

With $S=Q_O\mathsf E$ define $K=M_OS-SM_D$. For a strong-domain
field the exact identity is
\begin{equation}\label{eq:defect}
 Ku=(\PP_O-Q_O)\mathsf E M_Du
 +\PP_O\bigl[(1-\Delta_y)\mathsf E u
 -\mathsf E(1-\Delta_x)u-\mathsf E\nabla p_u\bigr].
\end{equation}
Using \eqref{eq:charts}, the principal second-order and pressure terms are
\begin{equation}\label{eq:principal-defect}
 \PP_O\left[-\chi R(GG^T-I)_{il}(\partial_{il}u)\circ F
 +\chi(G^T-R)(\nabla p_u)\circ F\right].
\end{equation}
The nonorthogonal pressure term is retained. Both coefficients have
supremum $Cr$ and Lipschitz norm $C$, including cutoff derivatives.
Their $C^\zeta$ norms are therefore at most $Cr^{1-\zeta}=Cr^{13/16}$.
Composition with the extended shear is bounded on $H^{-\eta,q}$,
uniformly in $r$, by duality and interpolation of its $L^{q'}$ and
$W^{1,q'}$ bounds. Multiplication by $C^\zeta$ functions is bounded
at order $-\eta$: the heat-kernel formula for
$[J_0^\eta,a]$ has near-zero integral
$\int_0^1t^{-1+(\zeta-\eta)/2}\,dt<\infty$, and duality gives the
negative-order bound. Together with \eqref{eq:negative-pressure},
\eqref{eq:principal-defect} is bounded by $Cr^{13/16}$ both as
$\Dom(M_D)\to H_q(O)$ and as
$\Dom(M_D^\theta)\to X_{-d,O}$.

The other Laplacian terms are
\[
 -(\Delta(\chi R))u\circ F
 -2\partial_a(\chi R)G_{ia}(\partial_i u)\circ F
 -\chi R(\Delta F_i)(\partial_i u)\circ F.
\]
The pressure remainder is $(\nabla\chi)p_u\circ F$, since
$\PP_O$ annihilates the flat gradient. They map
$\Dom(M_D^\theta)\to H_q(O)$ with bound $C_r$. The first term of
\eqref{eq:defect} has the same bound by
\eqref{eq:divergence-comparison} applied to $M_Du$ at order $-\eta$.

For $M_{i,\rho}=M_i+\rho$, $\rho\geq1$, we have proved both estimates
\begin{equation}\label{eq:two-norms}
 \norm{KM_{D,\rho}^{-1}}+
 \norm{M_{O,\rho}^{-d}KM_{D,\rho}^{d-1}}
 \leq Cr^{13/16}+C_r\rho^{-d}.
\end{equation}
For the lower-order term the first norm gains $\rho^{-d}$ from
$M_D^\theta M_{D,\rho}^{-1}$, the second from its left factor.
The principal terms are bounded on both fractional scales. Choose
$r$ first, then one $\rho(r)$ for all finitely many charts and both
dual exponents, so that the right side is at most
$a(r)=Cr^{13/16}$. The coefficient $C$ in this rate is uniform in $r$.

We give the root comparison quantitatively, since smallness on $H_q$
alone is insufficient for the trace reconstruction. Apply
\cite[Theorem 6.1, Lemma 6.2]{KKW} to
$\mathcal A=\diag(M_{D,\rho},M_{O,\rho})$ and the lower
off-diagonal perturbation $zK$. Its range condition is automatic
for the shifted invertible diagonal. The two smallness conditions
are exactly \eqref{eq:two-norms}. For $|z|a(r)<c$, the resolvent
series gives uniform $H^\infty$ bounds, and the small strong relative
bound gives uniform graph norms. Bounded imaginary powers and complex
interpolation \cite[Proposition 2.2]{KKW} give uniform square-root
norms on the shifted $V$ spaces. Explicitly, for $G=D,O$ set
\[
 \norm{v}_{V_{q,\rho}(G)}=\norm{M_{G,\rho}^{1/2}v}_{H_q(G)}
 \asymp\norm{v}_{V_q(G)}+\rho^{1/2}\norm{v}_{H_q(G)},\qquad\rho\geq1.
\]
The equivalence is uniform in $\rho$ by the isotropic calculus.

Conjugation of $\mathcal A$ by the triangular matrix with lower entry
$-zS$ gives the perturbation with lower entry $zK$. Its square root
has lower entry $zF_\rho$, where
\[
 F_\rho=M_{O,\rho}^{1/2}S-SM_{D,\rho}^{1/2},\qquad
 D_\rho=F_\rho M_{D,\rho}^{-1/2}.
\]
Taking $|z|$ comparable to $a(r)^{-1}$ in the uniform root bound gives
$\norm{D_\rho}_{H_q\to H_q}\leq Ca(r)$. The exact identity
\begin{equation}\label{eq:root-V}
 M_{O,\rho}^{1/2}D_\rho M_{D,\rho}^{-1/2}
 =KM_{D,\rho}^{-1}-D_\rho
\end{equation}
also gives $Ca(r)$ on the shifted $V_q$ norm. For $q=4$ the shifted
real interpolation norm satisfies, uniformly in $\rho$,
\[
 \norm{v}_{(H_4,V_{4,\rho})_{1/2,4}}
 \asymp\norm{v}_{\calE}+\rho^{1/4}\norm{v}_{H_4}.
\]
This follows by splitting the $K$-functional integral at
$\rho^{-1/2}$. Interpolation yields
\begin{equation}\label{eq:root-E}
 \norm{D_\rho v}_{\calE_O}
 \leq Ca(r)\norm{v}_{\calE}+C_r\norm{v}_{H_4},
\end{equation}
and the analogous $V_4$ estimate with a lower $H_4$ term.

Let $T_\rho=M_{O,\rho}^{1/2}SM_{D,\rho}^{-1/2}$.
Then $T-\mathsf E=D_\rho+\mathcal K_r$, with
$\mathcal K_r=(T-T_\rho)+(S-\mathsf E)$. To justify compactness at
the positive levels, keep the fixed shift and write
\begin{align}\label{eq:shift-difference}
 T-T_\rho={}&(M_O^{1/2}-M_{O,\rho}^{1/2})SM_D^{-1/2}\notag\\
 &+M_{O,\rho}^{1/2}S(M_D^{-1/2}-M_{D,\rho}^{-1/2}).
\end{align}
The first root difference equals
$-\rho(M_O^{1/2}+M_{O,\rho}^{1/2})^{-1}$ and gains one spatial
order. Since $S:V_q\to V_q$, its term maps $H_q$ into
$\Dom(M_O)$, compactly into $V_q(O)$. The inverse-root difference
in the second term maps $H_q$ to $\Dom(M_D^{3/2})$, compactly into
$\Dom(M_D)$. The strong-domain bound on $S$ then gives compactness
$H_q(D)\to V_q(O)$ for that term too. Together with
\eqref{eq:divergence-comparison} this makes $\mathcal K_r$ compact
on $H_q,V_q$ and, by interpolation, on $\calE$. These are statements
for the fixed selected $\rho$, with no uniform compactness assertion
as $\rho\to\infty$. The adjoint form argument in
Section~\ref{sec:energy-localizers} gives the consistent $Y_4$ action.

\subsection{Noise, drift, and reconstruction}\label{sec:reconstruction}

Write $B_{D,k}=\PP_D(b_k\cdot\nabla)$ and
$B_{j,k}=\beta_{j,k}\cdot\nabla$. For a no-slip solenoidal field,
\begin{equation}\label{eq:noise-pressure}
 B_{D,k}u=(b_k\cdot\nabla)u
 -(I-\PP_D)((u\cdot\nabla)b_k).
\end{equation}
The two unprojected derivatives have the same divergence and zero
normal component: this follows from $\divg b_k=\divg u=0$,
tangency of $b_k$, and the zero vector trace of $u$. Their difference
is solenoidal. The second term in \eqref{eq:noise-pressure} is bounded
on $H_4$; density gives the identity on $V_4$.

The exact commutator decomposition is
\begin{align}\label{eq:noise-commutator}
 T_jB_{D,k}-B_{j,k}T_j={}&(T_j-\mathsf E_j)B_{D,k}
 +(\mathsf E_jB_{D,k}-\beta_{j,k}\cdot\nabla\mathsf E_j)\notag\\
 &+\beta_{j,k}\cdot\nabla(\mathsf E_j-T_j).
\end{align}
The full constant derivative is used on the unprojected
$\mathsf E_ju$. The middle leading term is
$\chi_jR_j[(b_k\circ F_j)-G_j\beta_{j,k}]\cdot(\nabla u)\circ F_j$,
bounded by $Cr\norm{u}_{V_4}$. Cutoff, chart, and pressure terms
are bounded by $C_r\norm{u}_{H_4}$. The first and last terms use
the $H_4$ and $V_4$ bounds of $T_j-\mathsf E_j$, respectively.
Their compact remainders are compact $V_4\to H_4$, also after the
displayed derivatives. Summing the finite noise index in its
$\gamma$ norm gives a $C(a(r)+r)$ leading bound, modulo compact
$V_4\to H_4$ maps and lower $H_4$ norms.

The drift has the exact normalized identity
\begin{equation}\label{eq:drift-commutator}
 M_j^{-1/2}(M_{j,w}T_j-T_jM_{D,w})M_D^{-1/2}
 =K_jM_D^{-1}.
\end{equation}
It is small $V_4\to Y_4$, modulo a compact remainder: for the lower
term use $M_D^{-1/2}:V_4\to\Dom(M_D)$ followed by the compact
inclusion into $\Dom(M_D^\theta)$. Identity shifts cancel.

The partition and bounded overlap give
\begin{align}\label{eq:reconstruct-E}
 \norm{u}_{V_4}&\leq C\Bigl(\sum_j\norm{\mathsf E_ju}_{W^{1,4}}^4\Bigr)^{1/4}
 +C_r\norm{u}_{H_4},\\
 \norm{u}_{\calE}&\leq C\Bigl(\sum_j\norm{\mathsf E_ju}_{B^{1/2}_{4,4}}^4\Bigr)^{1/4}
 +C_r\norm{u}_{H_4}.\notag
\end{align}
For the second inequality use the $W^{1/2,4}=B^{1/2}_{4,4}$
difference integral. Pairs at distance less than $cr$ are reconstructed
from a uniformly bounded family of enlarged patches. The charts are
uniformly bi-Lipschitz and their Piola matrices uniformly Lipschitz.
Pairs at distance at least $cr$ cost only $C_r\norm{u}_{L^4}^4$.
This proves the uniform principal constant in \eqref{eq:reconstruct-E}.

Replacing $\mathsf E_j$ by $T_j$ costs at most
\begin{equation}\label{eq:count-absorption}
 CN^{1/4}a(r)\leq Cr^{-3/4}r^{13/16}=Cr^{1/16}\longrightarrow0.
\end{equation}
The leading noise and drift errors have the same bound. Choose $r$
small, then the common $\rho(r)$ of \eqref{eq:two-norms}. There
remain finitely many compact maps and lower $H_4$ norms. Ehrling's
inequality, with the continuous embedding into $L^2$, bounds them by
$\epsilon\norm{u}_{V_4}+C_{r,\epsilon}\norm{u}_{L^2}$, or at trace
level by $\epsilon\norm{u}_{\calE}+C_{r,\epsilon}\norm{u}_{L^2}$.
Here $\calE$ embeds compactly into $H_4$. Absorbing these terms
proves both reconstruction inequalities with $T_j$ and an $L^2$
remainder. For the continuous time norm use
$\sup_t\sum_j a_j(t)^4\leq\sum_j\sup_t a_j(t)^4$;
for the integral norm use Fubini. Thus the same reconstruction works
in $\calX_T$, including its actual continuous Besov trace.

\section{The physical linear equation}\label{sec:physical-linear}

We now prove Theorem~\ref{thm:linear}. Scale the physical noise to
$sB$, $0\leq s\leq1$. The covariances of every frozen model remain
between $\delta I$ and $\nu I$, and the vectors remain bounded.
The channel matrix limits, kernel and multiplier estimates, and
interior calculus are uniform on this compact parameter set. The
isotropic transfers do not depend on $s$. Thus the radius and common
shift in Section~\ref{sec:localization} can be fixed before choosing $s$.

For a solution already in $\calX_T$, apply the flat-model estimate to
$v_j=T_ju$. Its drift data are
$T_jf+\nu(A_{j,w}T_j-T_jA_{D,w})u$, and its noise is
\[
 sB_jv_j+(T_jsB_D-sB_jT_j)u+T_jg.
\]
The maps $T_j$ act consistently on $Y_4,H_4,\calE$ and preserve
predictability. Using \eqref{eq:noise-commutator}--\eqref{eq:count-absorption},
then the compact-error absorption, gives
\begin{equation}\label{eq:linear-apriori-lower}
 \norm{u}_{\calX_T}\leq C_r\bigl(\norm{h}_{L^4\calE}
 +\norm{f}_{L^4Y_4}+\norm{g}_{L^4\gamma(H_4)}\bigr)
 +C_r\norm{u}_{L^4(\Omega;C([0,T];L^2))}.
\end{equation}
The abbreviated data norms here and below are those of
Theorem~\ref{thm:linear}.

The remaining norm is controlled by physical energy. Orthogonality
of $\PP_D$ on $L^2$ and tangency and solenoidality of $b_k$ give
$(u,B_{D,k}u)_2=0$ and
\begin{equation}\label{eq:energy-coercivity}
 2\nu\norm{\nabla u}_2^2-s^2\sum_k\norm{B_{D,k}u}_2^2
 \geq2\delta\norm{\nabla u}_2^2.
\end{equation}
On bounded $D$, $V_2\hookrightarrow V_{4/3}$, so $Y_4\hookrightarrow V_2^*$;
also $H_4\hookrightarrow H_2$. The variational It\^o identity is
therefore applicable with the triple $V_2,H_2,V_2^*$. It can equally
be obtained by isotropic $L^2$ Stokes spectral projections, using
their uniform boundedness on $V_2,V_2^*$ and the drift and noise
$L^2$ bounds. The cross term with $g$ in the quadratic variation is
absorbed by $\delta\norm{\nabla u}_2^2$. The homogeneous-noise
martingale vanishes; the remaining martingale contains $(u,g_k)_2$.
The stochastic maximal inequality, followed by Young's inequality,
gives
\begin{equation}\label{eq:lower-energy}
 \norm{u}_{L^4(\Omega;CL^2)}\leq C_{T_0}
 \left[\norm{h}_{L^4(\Omega;L^2)}
 +T_0^{1/4}\bigl(\norm{f}_{L^4Y_4}+\norm{g}_{L^4\gamma(H_4)}\bigr)\right].
\end{equation}
The time factor follows from $L^4(0,T)\hookrightarrow L^2(0,T)$.
The estimate is uniform in $s$. Combining
\eqref{eq:linear-apriori-lower} and \eqref{eq:lower-energy} proves
\eqref{eq:linear-bound} as an a priori estimate. The choices were
$r$, then $\rho$, then the compact-error constants; none depends
on the solution, $T\leq T_0$, or $s$.

Existence follows by a finite continuation in $s$. At $s=0$, the
physical isotropic weak Stokes calculus and the additive theorem in
Section~\ref{sec:framework} give a causal solution operator $\mathcal S_0$.
If $\mathcal S_s$ exists, solve
\[
 u=\mathcal S_s(h,f,g+\Delta s\,B_Du)
\]
by contraction in predictable $\calX_T$. The constant is at most
$C\Delta s\norm{B_D}_{V_4\to\gamma(\R^K,H_4)}$, with the uniform
$C$ just proved. Choose one positive step making this less than
$1/2$; finitely many steps reach $s=1$. The deterministic and
stochastic integral estimates pass the limits to the equation in
$CY_4$. Energy proves linear uniqueness. Each step preserves
causality and restriction to common stopping times. This constructs
the operator asserted in Theorem~\ref{thm:linear}, rather than
inferring existence from an a priori bound.\qed

\section{The nonlinear equation}\label{sec:nonlinear}

The remaining argument uses only the physical spaces in
Section~\ref{sec:framework} and the constructed linear solution
operator. Define the convection term on $\calE$ by
\begin{equation}\label{eq:nonlinearity}
 \ip{F(u)}{\varphi}=\int_D(u\otimes u):\nabla\varphi,
 \qquad \varphi\in V_{4/3}.
\end{equation}
This is precisely $-\PP_D\divg(u\otimes u)$ in the weak formulation.

\subsection{Product bound and adapted local construction}\label{sec:cutoff}

We use $\calE\hookrightarrow L^8(D)$. A bounded extension of
$W^{1/2,4}(D)=B^{1/2}_{4,4}(D)$ to $\R^3$ is obtained from a
finite smooth atlas and even reflection. For cross-boundary pairs,
reflection only decreases the distance between the two reflected
points, so the fractional difference integral is bounded by the
half-space integral. Smooth bi-Lipschitz changes and cutoffs preserve
this bound. If $\Delta_jh$ is a dyadic decomposition of the extension,
Bernstein and H\"older give
\[
 \norm{h}_8\leq C\sum_{j\geq-1}2^{-j/8}
 (2^{j/2}\norm{\Delta_jh}_4)
 \leq C\norm{h}_{B^{1/2}_{4,4}},
\]
since $\sum_{j\geq-1}2^{-j/6}<\infty$. Consequently
\begin{equation}\label{eq:product}
 \norm{F(u)}_{Y_4}\leq C\norm{u}_{\calE}^2,\qquad
 \norm{F(u)-F(v)}_{Y_4}
 \leq C(\norm{u}_{\calE}+\norm{v}_{\calE})\norm{u-v}_{\calE}.
\end{equation}
This avoids any assertion that $\PP_D$ preserves zero vector trace
on $W^{1,4}$.

Choose a Lipschitz $\chi:[0,\infty)\to[0,1]$, equal to one on
$[0,1]$ and zero on $[2,\infty)$, and let
$F_R(u)=\chi(\norm{u}_{\calE}/R)F(u)$. It is bounded by $CR^2$
and globally Lipschitz with constant $L_R\leq CR$. Indeed, if
$\norm{u-v}_{\calE}<R$ and one cutoff is nonzero, both states have
norm less than $3R$; expand the cutoff difference and use
\eqref{eq:product}. If their distance is at least $R$, use the
$CR^2$ bound. If both cutoffs vanish there is nothing to prove.
For an adapted $\calE$-continuous process this is a predictable
state cutoff.

Let $\mathcal S_T$ denote the solution map in
Theorem~\ref{thm:linear}. In the closed Banach space of predictable
$\calX_T$ processes set $\Phi(u)=\mathcal S_T(h,F_R(u),0)$.
For deterministic $h,k\in\calE$ its estimates give
\[
 \norm{\Phi(u)-\Phi(v)}_{\calX_T}
 \leq C_L L_R T^{1/4}\norm{u-v}_{\calX_T}.
\]
Choose
\begin{equation}\label{eq:T-R}
 T_R=\min\{T_0,(2C_L\max(1,L_R))^{-4}\}.
\end{equation}
Contraction yields a predictable $U_R^h\in\calX_{T_R}$ and
\begin{equation}\label{eq:cutoff-lipschitz}
 \norm{U_R^h-U_R^k}_{\calX_{T_R}}\leq2C_L\norm{h-k}_{\calE}.
\end{equation}
The integral estimates pass the fixed-point limit to the truncated
equation in $CY_4$. For $R>\norm{h}_{\calE}$ set
\[
 \tau_R^h=\inf\{t\leq T_R:\norm{U_R^h(t)}_{\calE}\geq R\}\wedge T_R,
 \qquad \inf\varnothing=\infty.
\]
This stopping time is almost surely positive by continuity. On its
closed interval the cutoff is one, including at norm $R$, so
$u^h=U_R^h|_{[0,\tau_R^h]}$ solves the original equation. Its paths
have the desired Besov continuity and $L^4W^{1,4}_0$ regularity;
no $V_4$ value at the exit time is required.

\subsection{Uniqueness in the entire path class}\label{sec:uniqueness}

Let $u,v$ be any two adapted local solutions in
$CB^{1/2}_{4,4}\cap L^4W^{1,4}_0$, with the same initial datum and
Brownian motions, on lifetimes $\tau_u,\tau_v$. They have
$\calE$-continuous paths: the divergence and trace constraints hold
at almost every time from $V_4$, and then at every time by continuity
and closedness. No moments are assumed for these paths.
For deterministic $T$ stop at $\tau_u\wedge\tau_v\wedge T$ and
when the continuous running quantity
\[
 \norm{u(t)}_{\calE}+\norm{v(t)}_{\calE}
 +\int_0^t(\norm{u(s)}_{V_4}^4+\norm{v(s)}_{V_4}^4)\,ds
\]
first reaches $N$. Denote the resulting stopping time by $\sigma_N$.
These times increase to the common lifetime capped by $T$.

On the stopped interval $w=u-v$ belongs to $CH_2\cap L^2V_2$.
Convection lies in $L^2V_2^*$, since both states are bounded in
time in $L^4$, and $B_Dw\in L^2H_2$. The required second moments
are finite after this norm stopping. Apply It\^o's
formula to finite isotropic $L^2$ Stokes spectral projections until
$\sigma_N$. They converge on $V_2,V_2^*$ and are uniformly bounded
there; the drift pairings converge by the $L^2$ dual estimates,
quadratic variations by $L^2H_2$ convergence, and the martingales
by the isometry and maximal inequality. The $CH_2$ bound permits
uniform passage in time and then an almost sure subsequence.
This proves the variational identity until $\sigma_N$, as also
formulated in \cite[Lemma A.5]{AV}. It does not extend the stopped
state as a $V_4$-valued constant after $\sigma_N$.

By \eqref{eq:energy-coercivity} the noise correction leaves
$2\delta\norm{\nabla w}_2^2$. Also $(w,B_{D,k}w)_2=0$.
Convection by $u$ is skew. Integrating the derivative off $v$ in
the remaining term and using Young's inequality gives, for
$t\leq\sigma_N$,
\begin{equation}\label{eq:difference-energy}
 \norm{w(t)}_2^2+\delta\int_0^t\norm{\nabla w}_2^2\,ds
 \leq\norm{w(0)}_2^2+
 \delta^{-1}\int_0^t\norm{v}_\infty^2\norm{w}_2^2\,ds.
\end{equation}
The limiting martingale vanishes identically. Since
$W^{1,4}(D)\hookrightarrow L^\infty(D)$ and $L^4(0,T)\hookrightarrow L^2(0,T)$,
the Gronwall coefficient is finite pathwise. Equal initial data
give $w=0$ until $\sigma_N$. Let $N$ and then $T$ increase through
integers; path continuity includes the endpoints. This proves
uniqueness in the entire asserted class and identifies constructions
with different cutoff radii on their common lifetimes.

\subsection{Dependence after common stopping}\label{sec:dependence}

Fix $h$ and $R$ with $\norm{h}_{\calE}<R/4$. If smooth admissible
$h_n\to h$ in $\calE$, the same bound holds eventually for $h_n$.
Use one $T_R$, probability space, filtration, coefficient family and
Brownian coupling to construct $U=U_R^h$ and $U_n=U_R^{h_n}$. Set
\begin{align}\label{eq:common-stops}
 \tau&=\inf\{t\leq T_R:\norm{U(t)}_{\calE}\geq R/2\}\wedge T_R,\\
 \tau_n&=\inf\{t\leq T_R:\norm{U_n(t)}_{\calE}\geq R\}\wedge T_R,
 \qquad\sigma_n=\tau\wedge\tau_n.\notag
\end{align}
These stopping times are almost surely positive. Both restrictions
solve the original equation up to $\sigma_n$. For paths on this
interval define
\begin{equation}\label{eq:path-distance}
 d_\sigma(u,v)=\sup_{t\leq\sigma}\norm{u(t)-v(t)}_{\calE}
 +\left(\int_0^\sigma\norm{u-v}_{V_4}^4\,dt\right)^{1/4}.
\end{equation}
Restriction of \eqref{eq:cutoff-lipschitz} and Markov's inequality give
\begin{equation}\label{eq:dependence-bound}
 \PP\{d_{\sigma_n}(U_n,U)>\epsilon\}
 \leq C\epsilon^{-4}\norm{h_n-h}_{\calE}^4\longrightarrow0.
\end{equation}
If $\tau_n<\tau$, the perturbed path reached $R$ while the base
path had norm at most $R/2$. Its full-interval distance from $U$
was therefore at least $R/2$, and
\begin{equation}\label{eq:availability}
 \PP\{\sigma_n<\tau\}\leq CR^{-4}\norm{h_n-h}_{\calE}^4
 \longrightarrow0.
\end{equation}
The constants depend on the fixed coefficients, domain, $T_0,R$,
and not on $n$. On the event $\sigma_n=\tau$, both original
solutions reach the same base time $\tau$. Equations
\eqref{eq:dependence-bound}--\eqref{eq:availability} give continuity
in probability in the full path topology. The perturbed solution
reaches the base stopping time with probability tending to one.
No positivity of an infinite
infimum of exit times is needed. Together with whole-class uniqueness
and the local construction, this proves the main result.\qed

\paragraph{Use of artificial intelligence.}
OpenAI Codex was used extensively to develop the argument, inspect primary
sources, draft the manuscript, and perform internal mathematical and
editorial checks. Two additional Codex contexts read a fixed draft,
one as a reader and one as a referee, without the author checks or
each other's report. These same-provider AI reviews and the author
checks do not constitute human peer review or acceptance by the
problem repository.

The present revision also used Codex to derive the additional channel
regularity argument and recheck the corrected comparison theorem and its
dependent estimates. Those checks were performed in the same context as the
revision; no new independent referee report is claimed for this version.

\begin{thebibliography}{10}
\bibitem{AV}
A. Agresti and M. Veraar,
\emph{Nonlinear SPDEs and maximal regularity: an extended survey},
NoDEA Nonlinear Differential Equations Appl. \textbf{32} (2025), article 123.
\href{https://doi.org/10.1007/s00030-025-01090-2}{doi:10.1007/s00030-025-01090-2}.

\bibitem{AIM458}
AIM problem repository,
\emph{458. Local no-slip Navier--Stokes theory under non-small transport noise},
fixed problem statement.
\href{https://github.com/MColbrook/AIM/blob/a153f469ee0ea64e876b2ad97c79d252c19a3e22/problems/458-stochastic-navier-stokes-no-slip-local.md}{Problem statement}.

\bibitem{BG}
D. Breit and A. Gaudin,
\emph{Optimal regularity results for the Stokes--Dirichlet problem},
preprint, \href{https://arxiv.org/abs/2511.19091v1}{arXiv:2511.19091v1} (2025).

\bibitem{CM}
M. Costabel and A. McIntosh,
\emph{On Bogovski\u\i{} and regularized Poincar\'e integral operators for de Rham
complexes on Lipschitz domains}, author version, 29 December 2008.
\href{https://costabel.perso.math.cnrs.fr/publis/PoiBogdeRhv2.pdf}{Author manuscript}.

\bibitem{DK}
L. Deleaval and C. Kriegler,
\emph{Dimension free bounds for the vector-valued Hardy--Littlewood maximal operator},
preprint, \href{https://arxiv.org/abs/1703.08327v1}{arXiv:1703.08327v1} (2017).

\bibitem{GW}
M. Girardi and L. Weis,
\emph{Operator-valued Fourier multiplier theorems on $L^p(X)$ and geometry
of Banach spaces}, author version, 30 July 2002.
\href{https://people.math.sc.edu/girardi/papers/22.pdf}{Author manuscript}.

\bibitem{GoodairNavier}
D. Goodair,
\emph{Navier--Stokes Equations with Navier Boundary Conditions and Stochastic
Lie Transport: Well-Posedness and Inviscid Limit},
\href{https://arxiv.org/abs/2308.04290v2}{arXiv:2308.04290v2} (2023).

\bibitem{GoodairPlates}
D. Goodair,
\emph{Anisotropic Inviscid Limit for the Navier--Stokes Equations with
Transport Noise Between Two Plates},
preprint, \href{https://arxiv.org/abs/2603.13199v2}{arXiv:2603.13199v2} (2026).

\bibitem{KKW}
N. Kalton, P. Kunstmann, and L. Weis,
\emph{Perturbation and interpolation theorems for the $H^\infty$-calculus
with applications to differential operators},
Math. Ann. \textbf{336} (2006), 747--801.
\href{https://doi.org/10.1007/s00208-005-0742-3}{doi:10.1007/s00208-005-0742-3}.

\bibitem{KWErratum}
P. Kunstmann and L. Weis,
\emph{Erratum to: Perturbation and interpolation theorems for the
$H^\infty$-calculus with applications to differential operators},
Math. Ann. \textbf{357} (2013), 801--804.
\href{https://doi.org/10.1007/s00208-011-0768-7}{doi:10.1007/s00208-011-0768-7}.
Theorem~1.3 is the corrected comparison criterion used in
Section~\ref{sec:channel-consumer}; Theorem~2.1 supplies the physical
isotropic Stokes calculus.

\bibitem{PW}
J. Pr\"uss and M. Wilke,
\emph{On critical spaces for the Navier--Stokes equations},
\href{https://arxiv.org/abs/1703.08714v1}{arXiv:1703.08714v1} (2017).
\bibitem{Teschl}
G.~Teschl,
\emph{Mathematical Methods in Quantum Mechanics: With Applications to
Schr\"odinger Operators}, Graduate Studies in Mathematics 99,
American Mathematical Society, 2009. Author's online edition dated
12 February 2009, Theorem~2.13 and Section~3.1.
\url{https://www.mat.univie.ac.at/~gerald/ftp/book-schroe/schroe.pdf}.

\end{thebibliography}
\end{document}
